\chapter{Objets spectraux usuels.}
\VerdierPhysicalPageAnchor{0224}%
\VerdierSemanticAnchor{ch4d-4}
\begingroup
\small
\setlength{\parskip}{0pt}
\setlength{\emergencystretch}{1.5em}
\setlength{\abovedisplayskip}{4pt plus 2pt minus 2pt}%
\setlength{\belowdisplayskip}{4pt plus 2pt minus 2pt}%

\VerdierSemanticAnchor{ch4d-4-1}
\section{Premier objet spectral canonique.}

\VerdierSemanticAnchor{ch4d-4-1-1}
\noindent\textbf{4.1.1.}\quad
Rappelons que $\widetilde{\mathbb Z}$ désigne l'ensemble ordonné
$\mathbb Z\cup\{-\infty,+\infty\}$, ou encore la catégorie associée à cet
ensemble ordonné (chap. II, 4.3.1). Soient $\mathcal A$ une catégorie
additive et $Y\colon\widetilde{\mathbb Z}\longrightarrow
\operatorname{comp}(\mathcal A)$ un foncteur :

\VerdierSemanticAnchor{ch4d-4-1-1-1}
\begin{equation}
\begin{aligned}
q&\in\widetilde{\mathbb Z},
  &q&\longmapsto Y(q),\\
(p,q)&\in\mathcal{F}\ell_1(\widetilde{\mathbb Z}),
  &u(p,q)&\colon Y(p)\longrightarrow Y(q).
\end{aligned}
\tag{4.1.1.1}
\end{equation}

On a alors associé, au chapitre I, (3.3.6), au foncteur $Y$ un objet
spectral (de type $\widetilde{\mathbb Z}$) à valeurs dans $K(\mathcal A)$
(chap. II, 4.1.2 et 4.1.3) dont nous allons rappeler la construction.

Soit, tout d'abord, pour tout élément $(p,q)$ de
$\mathcal{F}\ell_1(\widetilde{\mathbb Z})$ (i.e.\ pour tout couple
$p,q\in\widetilde{\mathbb Z}$ tel que $p\leq q$), $\mathcal Y(p,q)$ le
cône du morphisme $u(p,q)$ (chap. I, 3.1.2) :

\VerdierSemanticAnchor{ch4d-4-1-1-2}
\begin{equation}
\begin{aligned}
\mathcal Y(p,q)^i&=Y(q)^i\oplus Y(p)^{i+1},\\
d_{\mathcal Y(p,q)}^i
&=
\begin{pmatrix}
d_{Y(q)}^i & u(p,q)^{i+1}\\
0 & -d_{Y(p)}^{i+1}
\end{pmatrix}.
\end{aligned}
\tag{4.1.1.2}
\end{equation}

Définissons alors, pour tout couple $(p,q)$ et $(p',q')$ d'éléments de
$\mathcal{F}\ell_1(\widetilde{\mathbb Z})$ tels que $p\leq p'$ et
$q\leq q'$ :

\VerdierSemanticAnchor{ch4d-4-1-1-3}
\begin{equation}
\Phi\bigl((p,q),(p',q')\bigr)\colon
\mathcal Y(p,q)\longrightarrow\mathcal Y(p',q'),
\tag{4.1.1.3}
\end{equation}
comme étant la classe à homotopie près du morphisme de complexes :

\VerdierSemanticAnchor{ch4d-4-1-1-4}
\begin{equation}
\Psi\bigl((p,q),(p',q')\bigr)\colon
\mathcal Y(p,q)\longrightarrow\mathcal Y(p',q'),
\tag{4.1.1.4}
\end{equation}
défini par les formules :

\VerdierSemanticAnchor{ch4d-4-1-1-5}
\begin{equation}
\Psi\bigl((p,q),(p',q')\bigr)^i
=
\begin{pmatrix}
u(q,q')^i & 0\\
0 & u(p,p')^{i+1}
\end{pmatrix}.
\tag{4.1.1.5}
\end{equation}
\endgroup
